% CLA Standard Attractor Benchmark Summary
% Generated 2026-07-07
% ASCII-only, pdflatex-safe

\documentclass[11pt]{article}
\usepackage[margin=1in]{geometry}
\usepackage{amsmath, amssymb}
\usepackage{booktabs}
\usepackage{hyperref}
\usepackage{listings}
\usepackage{xcolor}

\lstset{
  basicstyle=\ttfamily\small,
  frame=single,
  breaklines=true,
  showstringspaces=false,
}

\title{Chaos Language Algorithm: Standard Attractor Benchmark Summary}
\author{ZeroBot / ProtoCosmoBot}
\date{2026-07-07}

\begin{document}
\maketitle

\section{Overview}

This report summarises the Chaos Language Algorithm (CLA) pure-Python MVP
work on five standard chaotic dynamical systems. The CLA pipeline
comprises: trajectory generation, fixed-partition M1 symbolization,
n-gram chunk mining with non-overlapping rewrite and exact expansion,
deterministic context category proposal (exact or JS-divergence clustered),
greedy MDL acceptance loop, JSON persistence with edit-log replay, and
stable text rendering.

All code is in \texttt{projects/chaos-language-algorithm/repos/chaoslang}
(GitHub: \texttt{bgoertzel-sing/chaos-language-algorithm}), branch
\texttt{main}, latest commit \texttt{847390c}. All 64 unit tests pass.

\section{Systems Benchmark Suite}

Five standard chaotic systems were implemented as deterministic,
dependency-free Python generators:

\subsection{Logistic Map}

One-dimensional discrete map:
\begin{equation}
x_{t+1} = r \, x_t (1 - x_t), \quad r = 4.0, \; x_0 = 0.123
\end{equation}

Fully chaotic at $r=4$. Produces 4 unique M1 symbols (4 bins on $[0,1)$).
The CLA discovers 8 chunk rules including nested hierarchies
($N_0 \to N_3 \, N_9$, $N_0 \to N_4 \, N_5$). MDL score: 113.2 bits
(0.44 bits/symbol). Exact reconstruction verified.

\subsection{Lorenz-63}

Three-dimensional continuous-time flow:
\begin{align}
\dot{x} &= \sigma(y - x) \\
\dot{y} &= x(\rho - z) - y \\
\dot{z} &= xy - \beta z
\end{align}
with $\sigma = 10$, $\rho = 28$, $\beta = 8/3$, $dt = 0.01$, RK4 integration,
initial condition $(1, 1, 1)$.

The classic strange attractor. Produces 21 unique M1 symbols from the 3D
trajectory. CLA discovers 8 chunk rules with nested references. MDL score:
187.2 bits (0.73 bits/symbol). Exact reconstruction verified.

\subsection{R\"ossler}

Three-dimensional continuous-time flow:
\begin{align}
\dot{x} &= -y - z \\
\dot{y} &= x + ay \\
\dot{z} &= b + z(x - c)
\end{align}
with $a = 0.2$, $b = 0.2$, $c = 5.7$, $dt = 0.05$, RK4 integration,
initial condition $(0.1, 0, 0)$.

Simpler attractor structure than Lorenz-63 (single-band fold). Produces 24
unique M1 symbols. CLA discovers 8 chunk rules. MDL score: 189.2 bits
(0.74 bits/symbol). Exact reconstruction verified. Similar complexity to
Lorenz-63 as expected.

\subsection{Mackey-Glass}

Scalar delay-differential equation:
\begin{equation}
\frac{dx}{dt} = \beta \frac{x_{\tau}}{1 + x_{\tau}^n} - \gamma \, x
\end{equation}
with $\beta = 0.2$, $\gamma = 0.1$, $\tau = 17$, $n = 10$, $dt = 0.1$,
Euler method with discrete delay buffer, initial history $x = 0.5$.

Delayed-feedback chaotic system; the closest textbook analog to
OmegaSim-style delayed-coupling dynamics. Produces only 4 unique M1 symbols
(scalar trajectory, narrow dynamic range with these parameters). CLA
discovers 6 chunk rules including nested hierarchy
($N_0 \to N_1 \to N_5$). Best compression ratio: 0.17 bits/symbol. MDL
score: 44.4 bits. Exact reconstruction verified.

\subsection{Lorenz-96}

$N$-dimensional continuous-time flow:
\begin{equation}
\dot{x}_i = (x_{i+1} - x_{i-2}) \, x_{i-1} - x_i + F, \quad
i = 0, \ldots, N-1 \pmod{N}
\end{equation}
with $F = 8.0$, $dt = 0.01$, RK4 integration. Dimension inferred from
initial condition. Default: $N = 5$ with initial $[8.0, 8.0, 8.0, 8.0, 8.01]$.

Chaotic for $F \geq 4$ and $N \geq 4$. This is the primary scalability
benchmark because dimension is freely configurable.

\section{Dimensionality Scaling Results (Lorenz-96)}

Benchmark runs at dimensions $N \in \{5, 8, 12, 16, 20, 24\}$ with
4 bins per axis, 256 post-discard steps, 8 MDL iterations:

\begin{table}[h]
\centering
\begin{tabular}{ccccc}
\toprule
Dim & Unique Symbols & Rules & MDL Score (bits) & Bits/Symbol \\
\midrule
5  & 56  & 8 & 234.2 & 0.91 \\
8  & 78  & 8 & 229.2 & 0.90 \\
12 & 107 & 6 & 241.4 & 0.94 \\
16 & 121 & 3 & 248.7 & 0.97 \\
20 & 127 & 6 & 251.4 & 0.98 \\
24 & 148 & 5 & 251.5 & 0.98 \\
\bottomrule
\end{tabular}
\caption{Lorenz-96 dimensionality scaling. At dim $\geq 16$, the symbol
stream approaches near-incompressibility with 4 bins and 256 steps:
most symbols are unique. CLA maintains exact reconstruction at all
tested dimensions.}
\end{table}

Key observations:
\begin{itemize}
\item Unique symbols grow rapidly with dimension: $4^N$ possible symbols,
so at $N=16$ there are $\sim 4.3 \times 10^9$ possible symbols versus
256 actual observations. Most are unique.
\item MDL score approaches the raw-entropy ceiling ($\log_2 256 = 8$ bits/symbol
is the trivial upper bound; actual unique-symbol entropy approaches
$\sim 8$ bits/symbol at high $N$).
\item Chunk rule count drops at high dimension because few repeated n-grams
exist in 256 steps. Longer trajectories or coarser binning would help.
\item At dim $= 20$, the experiment notes state this is comparable to
OmegaSim starter-vector dimensionality ($\sim 17$ state variables per hive,
$\sim 20$--$30$ when including cross-hive coupling fields).
\item At dim $= 24$, 58\% of symbols are unique, yet CLA still finds
5 repeatable chunk rules and maintains exact reconstruction.
\end{itemize}

\section{CLA Pipeline Architecture}

\subsection{Trajectory Generation}

Five deterministic generators in \texttt{benchmarks/attractors.py}:
\texttt{logistic\_map}, \texttt{lorenz63}, \texttt{rossler},
\texttt{mackey\_glass}, \texttt{lorenz96}. Continuous systems use
fixed-step RK4; Mackey-Glass uses Euler with discrete delay buffer;
logistic map is a direct discrete iteration. All are seed-deterministic
(no RNG).

\subsection{M1 Symbolization}

Fixed rectangular partition: for each trajectory point $\mathbf{x} \in
\mathbb{R}^d$, each axis $i$ is binned into $B$ equal-width bins over
$[\min_i, \max_i]$. The symbol is a concatenated label:
\texttt{m1:d0\_2|d1\_0|d2\_3} for a 3D point. Bounds are inferred once
from the full trajectory (or supplied externally for cross-trajectory
comparison). Values outside supplied bounds are clamped to edge bins.

\subsection{Chunk Mining}

N-gram pattern miner (\texttt{mining/ngram.py}): scans for repeated
n-grams of length $n \in [n_{\min}, n_{\max}]$ with minimum use count.
Non-overlapping rewrite replaces occurrences with a new nonterminal
token. Pruning removes productions whose nonterminal is no longer
referenced after rewrite (fixes dangling-nonterminal bug).

\subsection{Context Categories}

Two category induction methods:
\begin{enumerate}
\item \textbf{Exact} (\texttt{ContextCategoryInducer}): groups tokens
by exact context signature (left/right neighbour distribution).
Category occurrences preserve the observed member for exact reconstruction.
\item \textbf{JS-divergence} (\texttt{JSCategoryInducer}): clusters tokens
whose context histograms have Jensen-Shannon divergence below a
threshold. Implemented in \texttt{categorization/js.py} with greedy
sorted-key clustering against centroid. Integrated into category proposal
generation as an alternative to exact grouping.
\end{enumerate}

\subsection{MDL Scoring}

Two-part MDL (\texttt{scoring/mdl.py}): model cost = sum of rule
overhead ($0.8$ bits/production) + production RHS length + category
overhead ($14.0$ bits/category) + member set size + edit-log cost
($0.1$ bits/edit). Data cost = token count for base tokens +
$1 - \log_2 p(\text{member}|\text{category})$ for category occurrences.
Greedy acceptance: a chunk/category proposal is accepted iff total MDL
decreases.

\subsection{Persistence}

JSON state serialization (\texttt{persistence.py}): round-trips
\texttt{GrammarState} to/from JSON with sorted keys for determinism.
Edit-log replay reconstructs state from initial symbols + edit sequence.
\texttt{grammar\_to\_text} provides stable human-readable grammar
rendering. All round-trips preserve exact reconstruction.

\section{Common Benchmark Parameters}

All experiments use:
\begin{itemize}
\item \texttt{steps = 256} (post-discard)
\item \texttt{discard = 64} (transient removal)
\item \texttt{bins = 4} (M1 partition granularity per axis)
\item \texttt{iterations = 8} (MDL loop iterations)
\item \texttt{seed = 0} (for CLA's internal RNG used in tie-breaking)
\end{itemize}

\section{OmegaSim Dimensionality Context}

The current OmegaSim A6 single-hive model has 17 state variables per
hive: 4 motivation, 4 fatigue, 4 threshold, plus scalars (artifact,
provenance\_debt, risk, prediction\_error, last\_action). A three-hive
A8 configuration would have $\sim 51$ state variables plus inter-hive
coupling fields.

The CLA Lorenz-96 scaling experiments at dim $= 20$--$24$ are directly
comparable to single-hive OmegaSim dimensionality. At these dimensions CLA
maintains exact reconstruction but approaches raw-entropy compression
limits with 4 bins and 256 steps.

The next step is selecting toy attractor systems with dimensionality
$\sim 17$--$30$ to bridge from textbook benchmarks to OmegaSim-scale
traces, and testing whether CLA can recover known grammatical structure
at that scale with longer trajectories and/or coarser binning.

\section{Summary Table}

\begin{table}[h]
\centering
\begin{tabular}{lcccc}
\toprule
System & Dim & Unique Symbols & Rules & Bits/Symbol \\
\midrule
Logistic Map   & 1  & 4   & 8 & 0.44 \\
Mackey-Glass   & 1  & 4   & 6 & 0.17 \\
Lorenz-63      & 3  & 21  & 8 & 0.73 \\
R\"ossler      & 3  & 24  & 8 & 0.74 \\
Lorenz-96 (5)  & 5  & 56  & 8 & 0.91 \\
Lorenz-96 (8)  & 8  & 78  & 8 & 0.90 \\
Lorenz-96 (12) & 12 & 107 & 6 & 0.94 \\
Lorenz-96 (16) & 16 & 121 & 3 & 0.97 \\
Lorenz-96 (20) & 20 & 127 & 6 & 0.98 \\
Lorenz-96 (24) & 24 & 148 & 5 & 0.98 \\
\bottomrule
\end{tabular}
\caption{All 5 standard attractor systems with Lorenz-96 dimensionality
scaling. All runs: 256 steps, 4 bins, 8 iterations, exact reconstruction
verified.}
\end{table}

\section{Repository State}

\begin{itemize}
\item Repository: \texttt{bgoertzel-sing/chaos-language-algorithm}
\item Branch: \texttt{main}
\item Latest commit: \texttt{847390c} (persistence complete)
\item Tests: 64 passing (Python 3.9 on MacBook)
\item Dependencies: pure Python stdlib (no numpy, scipy, or torch)
\item Sprint tasks 1--5 complete: symbolic MVP, M1 attractor controls,
benchmark experiments, JS-divergence clustering, persistence.
\end{itemize}

\end{document}
